\documentclass[10pt,a4paper]{amsart}
\usepackage[margin=19mm]{geometry}
\usepackage{amsmath,amssymb,mathtools}
\usepackage[scr=rsfs,cal=euler]{mathalfa}
\usepackage{xfrac}
\usepackage[unicode,hidelinks,bookmarks=false]{hyperref}
\newcommand{\ie}{i.e.\ }
\newcommand{\cf}{cf.\ }
\newcommand{\resp}{respectively\ }
\newcommand{\Subsection}{\subsection}
\providecommand{\nobreakdash}{\nobreak\hbox{-}}
\setlength{\emergencystretch}{2em}
\multlinegap0pt
\usepackage{amssymb}
\usepackage{algorithm}\usepackage{algpseudocode}
\usepackage{tcolorbox}
\usepackage{xcolor}
\newtheorem{theorem}{Theorem}
\newtheorem{proposition}{Proposition}
\title{Hybrid Energy-Aware Reward Shaping:\\A Unified Lightweight Physics-Guided Methodology for Policy Optimization}
\author{Qijun Liao, Jue Yang, Yiting Kang, Xinxin Zhao, Yong Zhang,, Mingan Zhao}
\date{Source: arXiv:2603.11600v2 (CC BY 4.0)}
\begin{document}
\maketitle

\noindent\emph{Source and licence.} Hybrid Energy-Aware Reward Shaping:\\A Unified Lightweight Physics-Guided Methodology for Policy Optimization. Version arXiv:2603.11600v2 (CC BY 4.0), distributed by arXiv under the Creative Commons Attribution 4.0 International licence, http://creativecommons.org/licenses/by/4.0/, as recorded in the arXiv licence metadata for this exact version and shown on the abstract page https://arxiv.org/abs/2603.11600v2. Full licence notice: \texttt{licenses/arxiv-2603.11600-CC-BY-4.0.txt}.\par\smallskip

\noindent\textit{Editorial scope.} Counted unit: the article's proposition prop:dual\_necessity (source0.tex lines 434--447). The article's complete original proof of it is delivered with it (source0.tex lines 449--485, one \texttt{proof} environment of 37 lines). No mathematics is written, repaired or re-derived: the statement and the proof are the article's own lines, in the article's own notation, and no retained line is substituted or re-flowed.  Retained uncounted as the source-local context the proof needs: lines 291--312.  Nothing was omitted from any retained span, so the package carries no omission record.  No retained line contains a citation, so the wrapper carries no bibliography and no bibliography entry.  No retained line contains a figure environment, an image-inclusion command or drawing code, so no image is delivered and the package declares no asset dependency.  Editorial formatting only, in addition: an AMS \texttt{amsart} preamble that restates the article's own declarations and macros used by the retained lines, namely the environments \texttt{theorem}, \texttt{proposition} on separate counters, together with the standard packages those lines need.  The retained environments print in this document as Theorem 1, Proposition 1 in order of appearance, whereas the article prints the same items with the numbering of its own full document.  The complete native source of the article as distributed in the arXiv e-print for this exact version is bundled unchanged as \texttt{sources/196226/source0.tex}.
	\begin{theorem}[Energy-Based Gradient Enrichment Under Mechanical Stability]
		\label{thm:energy_convergence}
		Consider a mechanical system with generalized coordinates $q\in\mathbb{R}^n$, velocities $\dot{q}\in\mathbb{R}^n$, and total energy $E(q,\dot{q}) = T(\dot{q}) + U(q)$. Define $\Phi_{\text{energy}}(s) \triangleq -E(q(s),\dot{q}(s))$.
		
		If the mechanical stability condition
		\begin{equation}
			\frac{\partial^2 E}{\partial q^2}\bigg|_{q \in \mathcal{S}_{\text{op}}} \succ 0
			\label{eq:mechanical_stability}
		\end{equation}
		holds over the operationally relevant region $\mathcal{S}_{\text{op}}\subseteq\mathcal{S}$, then the shaped reward $\tilde{R}(s,a,s') = R(s,a,s') + \gamma\Phi_{\text{energy}}(s') - \Phi_{\text{energy}}(s)$ induces a policy gradient decomposition
		\begin{align}
			\nabla_\theta J_{\text{shaped}}(\theta) &= \nabla_\theta J_{\text{original}}(\theta) \notag \\
			&\quad + \alpha_{\text{energy}} \mathbb{E}_{\tau\sim\pi_\theta} \bigg[ \sum_{t=0}^{\infty} \gamma^t \nabla_\theta \log \pi_\theta(a_t|s_t) \cdot \Delta\Phi_{\text{energy},t} \bigg]
			\label{eq:gradient_decomposition}
		\end{align}
		where $\Delta\Phi_{\text{energy},t} = \gamma\Phi_{\text{energy}}(s_{t+1}) - \Phi_{\text{energy}}(s_t)$. Moreover, condition~\eqref{eq:mechanical_stability} guarantees $\|\nabla_q E(q)\| \geq \varepsilon_{\mathcal{K}} > 0$ throughout any compact $\mathcal{K}\subset\mathcal{S}_{\text{op}}\setminus\mathcal{B}_r(q^*)$, so that the gradient informativeness ratio
		\begin{equation}
			\rho_{\text{info}} \triangleq \frac{\|\nabla_q E\|_{\text{energy-informative}}}{\|\nabla_s R\|_{\text{task-sparse}}}
			\label{eq:gradient_ratio}
		\end{equation}
		satisfies $\rho_{\text{info}} \gg 1$ in sparse-reward locomotion tasks where condition~\eqref{eq:mechanical_stability} holds.
	\end{theorem}

	\begin{proposition}[Dual-Potential Decomposition Necessity]
		\label{prop:dual_necessity}
		For control tasks where task objectives (e.g., reaching targets, tracking trajectories) and energy efficiency are not naturally aligned, a single potential function $\Phi_{\text{single}}(s)$ cannot simultaneously satisfy:
		\begin{enumerate}
			\item Task directivity: $\nabla_s \Phi_{\text{single}}$ points toward task-relevant states
			\item Energy awareness: $\Phi_{\text{single}}$ encodes mechanical energy structure
		\end{enumerate}
		
		The dual decomposition $\Phi(s) = \alpha_{\text{task}}\Phi_{\text{task}}(s) + \alpha_{\text{energy}}\Phi_{\text{energy}}(s)$ resolves this conflict by enabling independent tuning of:
		\begin{equation}
			\nabla_\theta J_{\text{shaped}} = \nabla_\theta J_{\text{original}} + \underbrace{\alpha_{\text{task}} \nabla_\theta J_{\Phi_{\text{task}}}}_{\text{task guidance}} + \underbrace{\alpha_{\text{energy}} \nabla_\theta J_{\Phi_{\text{energy}}}}_{\text{energy guidance}}
		\end{equation}
		where $\alpha_{\text{task}}$ and $\alpha_{\text{energy}}$ control the balance between exploration efficiency and physical plausibility.
	\end{proposition}

	\begin{proof}
		The proof proceeds by contradiction, demonstrating incompatibility of dual requirements.
		
		\textbf{Assumption}: Suppose a single potential function $\Phi_{\text{single}}:\mathcal{S}\to\mathbb{R}$ can achieve both task directivity and energy awareness. Then for navigation tasks:
		
		Task requirement: To guide toward goal state $s_{\text{goal}}$, the potential must satisfy:
		\begin{equation}
			\Phi_{\text{single}}(s) \approx -\|s - s_{\text{goal}}\|^2 + C_1
			\label{eq:task_req}
		\end{equation}
		ensuring $\nabla_s \Phi_{\text{single}}$ points toward $s_{\text{goal}}$.
		
		Energy requirement: To encode mechanical structure, the potential must satisfy:
		\begin{equation}
			\Phi_{\text{single}}(s) = -E(q(s),\dot{q}(s)) + C_2 = -\left[T(\dot{q}) + U(q)\right] + C_2
			\label{eq:energy_req}
		\end{equation}
		ensuring energy minimization properties (Theorem~\ref{thm:energy_convergence}).
		
		Contradiction: These requirements conflict when the minimum-energy path differs from the shortest geometric path. Specifically:
		
		\begin{itemize}
			\item Geometric shortest path: For point-to-point navigation, Eq.~\eqref{eq:task_req} induces straight-line trajectories minimizing $\|s-s_{\text{goal}}\|$.
			
			\item Energy-efficient path: For systems with kinetic energy $T=\frac{1}{2}\dot{q}^\top M\dot{q}$, Eq.~\eqref{eq:energy_req} favors smooth, low-acceleration trajectories satisfying $\ddot{q}^\top M\ddot{q}$ minimization (principle of least action).
			
			\item Incompatibility: When obstacles or dynamics constraints exist, the energy-efficient path may detour to maintain smooth motion, conflicting with the straight-line preference of Eq.~\eqref{eq:task_req}.
		\end{itemize}
		
		Resolution through decomposition: 
		\begin{itemize}
			\item Early training ($\alpha_{\text{task}} \gg \alpha_{\text{energy}}$): Prioritize task completion to establish basic competence, accepting energy-inefficient exploration.
			\item Late training ($\alpha_{\text{task}} \approx \alpha_{\text{energy}}$): Refine policies toward energy-efficient solutions while maintaining task success.
		\end{itemize}
		
		This staged optimization is impossible with $\Phi_{\text{single}}$ because it forces a fixed trade-off encoded in the function itself, whereas the decomposition enables dynamic balancing through $\alpha_{\text{task}}, \alpha_{\text{energy}}$ tuning.
	\end{proof}

\end{document}
